\scanpage{24}
\mspage{7}
de $\dot X\subset\widetilde X$ est déterminé par la connaissance de
\[
\widetilde X'_1\cap\dot X=\dot X'_1.
\]

\uline{Corollaire} Les parties \inserted{fermées} $X'$ de $X$ compatibles avec
le germe de structure cônique, sont en
corr. $1$-$1$ avec les couples $(\widetilde X'\subset\widetilde X,Y'\subset Y)$
de parties fermées de $\widetilde X$ et de $Y$, telles que

a) $\widetilde X'$ soit compatible avec le \inserted{germe de} structure cylindrique
\inserted{de $\widetilde X,\dot X$}

b) $f:\dot X\to Y$ applique \deleted{$\unsure{\widetilde X'}$}
$\dot X'\overset{\mathrm{def}}=\widetilde X'\cap\dot X$ dans $Y'$.

On aura
\[
X'=\operatorname{Im}\underbrace{(\widetilde X'\hookrightarrow\widetilde X\to X)}
_{\overline{X'\setminus\underbrace{Y\cap X'}_{Y'}}}\cup Y'.
\]
Pour que $X'$ soit \deleted{\uline{strict.} compatible avec}
\inserted{cônique \uline{\unsure{strict}}},
\deleted{\inserted{\unsure{purement}} la structure cônique}
il f. et s. que $Y'=\operatorname{Im}(\dot X'\to Y)$

Soit maintenant $X_0\subset X$ une partie
bordante de $X$, et $\mathcal E$ une \inserted{germe de} structure cylindrique
de $X_0$ dans $X$. On dit que $\mathcal E$ est compatible
avec le germe de structure cônique $\mathcal T$
de $X$ autour de $Y$, si $\exists$ représentants
$E\inserted{\simeq X_0\times I}$ voisinage cylindrique de $X_0$, \unsure{et} $T$ voisinage
cônique de $Y$ (avec \inserted{sa} structure cônique), tels que

\margin{au lieu de ``compatible'' il \unsure{est} plus \unsure{parlant} de
dire ``\uline{transversal}''.}

a) $\overline T\cap E$, $\dot T\cap E$, et $Y\cap E$ \deleted{\unsure{compat.}}
\inserted{soient} des parties cylindriques de
\deleted{$\unsure{\overline E}\cap\overline T$ soit stable par les deux types d'homothéties
$\unsure{\overline x}\mapsto\unsure{\lambda x}_{X_0}$ et $t\mapsto\lambda\times_Y t$
qui \unsure{agissent} \unsure{de part et d'autre}}
$E$ i.e.

\margin{\unsure{donc} stable par $E$-homothéties par $\lambda\in[0,1]$}

\deleted{
\diagram{
\begin{tikzcd}[column sep=large,row sep=large]
E \arrow[r,phantom,"\simeq"] &
X_0\times I \\
\overline T\cap E \arrow[u,phantom,"\cup" description]
\arrow[r,phantom,"\simeq"] &
\overline T_0\times I \arrow[u,phantom,"\cup" description] \arrow[d] &
\overline T_0=X_0\cap\overline T \\
Y\cap E \arrow[r,phantom,"\simeq"] &
Y_0\times I &
Y_0=X_0\cap Y
\end{tikzcd}}
$\deleted{\unsure{E\cap T}}$
}
\margin{23}
